%%%%%%%%%%  scan idx 201  |  folio 188  %%%%%%%%%%
isomorphisme fonctoriel en $E$ :
\[
 \Hom(E,F)\simeq\Hom(E,QF).
\]
Cela prouve l'existence de $Q$; l'unicité est évidente puisqu'il s'agit d'un adjoint.

\textbf{Lemme 3.3.} \begin{enumerate}[label=\alph*)]
\item Le cohéreur est un foncteur additif transformant injectifs en injectifs.
\item Notons
\[
 \mathbb RQ:\mathrm D^+(S)\longrightarrow \mathrm D^+(\Qcoh(S))
\]
le dérivé droit du cohéreur. Pour $E\in\Ob\mathrm D(\Qcoh(S))$ et $F\in\Ob\mathrm D^+(S)$, il existe des isomorphismes fonctoriels canoniques
\begin{align}
 \mathbb R\Hom_{\Mod}(E,F)&\simeq \mathbb R\Hom_{\Qcoh}(E,\mathbb RQ(F)) \qquad\text{\footnotemark}\tag{i}\\
 \Hom_{\Mod}(E,F)&\simeq \Hom_{\Qcoh}(E,\mathbb RQ(F)). \tag{ii}
\end{align}\footnotetext{L'indice $\Mod$ (resp. $\Qcoh$) est une abréviation pour $\mathrm D(S)$ (resp. $\mathrm D(\Qcoh(S))$).}
\end{enumerate}

\emph{Preuve.} L'assertion a) est conséquence immédiate de la définition. Pour b), il suffit d'établir l'isomorphisme (i), car (ii) s'en déduira par passage à la cohomologie. On peut, pour cela, supposer $F$ à composantes injectives; $QF$ est alors à composantes injectives d'après a), et l'isomorphisme (i) provient simplement de la définition de $Q$.

\textbf{Lemme 3.4.} Soit $f:X\longrightarrow Y$ un morphisme de schémas quasi-compacts et quasi-séparés. Pour $E\in\Ob\Mod(X)$, il existe un isomorphisme fonctoriel canonique
\[
 Qf_*(E)\simeq f_*Q(E).
\]

